The proof combines existing non-asymptotic tools with RR-specific algebra and
 the deterministic-start transfer developed here. The main proof blocks are:

\begin{table}[htbp]
\centering
\small
\caption{Main proof blocks.}
\begin{tabular}{|p{0.19\textwidth}|p{0.36\textwidth}|p{0.34\textwidth}|}
\hline
\textbf{Block} & \textbf{Role in the proof} & \textbf{Status} \\
\hline
RR weight algebra &
Controls the deterministic kernels $\mathcal{Q}_{l}^{\mathrm{RR}}$ and
$Q_{l;n_{0},n}^{\mathrm{RR}}$ and their variation. &
Derived in \cref{sec:zeroth_order_rr,sec:pr_weights}. \\
\hline
Poisson and martingale reduction &
Turns the Markovian depth-zero weighted sum into a martingale plus a bounded
Poisson remainder. &
Adapted from the Samsonov et al. framework and reproved for RR weights. \\
\hline
Stationary misadjustment &
Controls the non-martingale RR remainder
$J^{(1)}+J^{(2)}+H^{(2)}$ under the stationary augmented chain. &
Uses Levin et al. inputs; the stationary RR assembly is carried out here. \\
\hline
Burn-in transfer &
Transfers the stationary theorem to finite starts by controlling the
deterministic transient, random initial product, and startup discrepancy. &
Developed in \cref{sec:burn_in_transfer}. \\
\hline
\end{tabular}
\end{table}

The theorem-level outputs should be read in the following order:

\begin{table}[htbp]
\centering
\small
\caption{Theorem-level outputs and labels.}
\begin{tabular}{|p{0.25\textwidth}|p{0.23\textwidth}|p{0.41\textwidth}|}
\hline
\textbf{Output} & \textbf{Label} & \textbf{Object controlled} \\
\hline
Stationary augmented-chain assembly &
\cref{thm:RR-BE} &
$S_{n,\mathrm{stat}}^{\mathrm{RR}}(u)$ with the finite-window variance
$\sigma_{n}^{\mathrm{RR}}(u)$. \\
\hline
Stationary balanced triangular-array corollary &
\cref{cor:RR-BE-working} &
$S_{n,\mathrm{stat}}^{\mathrm{RR}}(u)$ at $\alpha_n=cn^{-1/2}$, with either
finite-window normalization. \\
\hline
Deterministic-start finite-window transfer &
\cref{thm:burn-RR-BE-master} &
$\Xi_{n,n_{0}}^{\mathrm{bRR}}(u)$ after burn-in, still normalized by
$\sigma_{n,n_{0}}^{\mathrm{bRR}}(u)$. \\
\hline
Deterministic-start balanced theorem &
\cref{thm:burn-final-balanced,cor:burn-sqrt-n-transfer} &
$\Xi_{n,n_{0}}^{\mathrm{asy,RR}}(u)$ and then the final $\sqrt n$ statistic
$\Xi_{n,n_{0}}^{\mathrm{n,RR}}(u)$. \\
\hline
\end{tabular}
\end{table}

The following notation is used throughout the later chapters:

\begin{table}[htbp]
\centering
\small
\caption{Notation used in the later chapters.}
\begin{tabular}{|p{0.26\textwidth}|p{0.64\textwidth}|}
\hline
\textbf{Notation} & \textbf{Meaning} \\
\hline
$\alpha,2\alpha$ & The two constant step sizes used by the two-level RR estimator. \\
\hline
$n_0$ and $m=n-n_0$ & Burn-in length and effective averaging window. \\
\hline
$(Z_k)_{k\ge0}$ and $\xi$ & Base Markov chain and initial law
$\xi=\mathcal{L}(Z_0)$. The recursion uses observations $Z_k$, $k\ge1$;
stationary covariance formulas may use $Z_0$ by stationarity. \\
\hline
$\mathcal{F}_k$ & Natural filtration $\sigma(Z_0,\ldots,Z_k)$. Local displays
that begin at $Z_1$ use the same filtration with the harmless extra variable
$Z_0$. \\
\hline
$(Z_k^{\mathrm{stat}})_{k\in\mathbb Z}$ & Two-sided stationary copy of the
base chain. Under the stationary augmented-chain convention the superscript is
usually omitted. \\
\hline
$B_w=I-w\overline A$ & Deterministic linearized contraction at step size $w$. \\
\hline
$\mathcal{Q}_l^{\mathrm{RR}}$ & Full-window RR PR weight in the stationary
$n_0=0$ chapter. \\
\hline
$Q_{l;n_0,n}^{\mathrm{RR}}$ or $Q_l^{\mathrm{bRR}}$ & Burned-in RR PR weight
for the window $k=n_0,\ldots,n-1$. \\
\hline
$S_{n,\mathrm{stat}}^{\mathrm{RR}}(u)$ & Stationary augmented-chain scalar statistic. \\
\hline
$T_{n,n_0}^{\mathrm{RR}}(u)$ & Finite-start burned-in scalar statistic with
$\sqrt m$ normalization. \\
\hline
$\Xi_{n,n_0}^{\mathrm{bRR}}(u)$ & Finite-window normalized burned-in statistic. \\
\hline
$\sigma(u)^2=u^{\mathsf T}\Sigma_\infty u$ & Asymptotic scalar variance target. \\
\hline
$\operatorname{polylog}(n)$ & A generic fixed power of $\log n$, possibly
changing from line to line. It never hides polynomial dependence on $n$ or any
dependence on $\alpha$, $p$, or $q$ that is rate-relevant in the surrounding
display. \\
\hline
\end{tabular}
\end{table}
